\section{GRPO: Group Relative Policy Optimization}

\begin{frame}{GRPO Overview}
    \begin{itemize}\itemsep0.35em
        \item \textbf{GRPO}: Group Relative Policy Optimization, a PPO-style clipped update with the critic removed.
        \item PPO needs a learned value network to say how good a response was \emph{expected} to be. GRPO gets that baseline for free: sample several responses to the same prompt and score each against the group's mean reward.
        \item So one fewer network to train and hold in memory, and a baseline that cannot lag behind the policy because it is recomputed from scratch at every step.
        \item Introduced in DeepSeekMath (Shao et al., 2024), then used for the RL stages of DeepSeek-V2, DeepSeek-V3, and DeepSeek-R1.
    \end{itemize}
\end{frame}

\begin{frame}[allowframebreaks]{GRPO Objective}
    \begin{itemize}
        \item GRPO avoids the need for a separate critic model by estimating the baseline from group scores.
        \item For each question $q$, sample a group of outputs $\{o_1, o_2, \ldots, o_G\}$ from the old policy $\pi_{\theta_{\text{old}}}$.
        \item The policy model $\pi_\theta$ is optimized by maximizing:
        \begin{equation}
        \begin{split}
            J_{\text{GRPO}}(\theta) =\; & \mathbb{E}_{q \sim P(Q),\, \{o_i\}_{i=1}^G \sim \pi_{\theta_{\text{old}}}(O|q)} \Bigg[
            \frac{1}{G} \sum_{i=1}^G \frac{1}{|o_i|} \sum_{t=1}^{|o_i|} \Bigg\{
                \min \Bigg(
                    w_{i,t}(\theta)\, \hat{A}_{i,t}, \\
                & \qquad\qquad
                    \text{clip}\big(w_{i,t}(\theta),\, 1-\epsilon,\, 1+\epsilon\big) \hat{A}_{i,t}
                \Bigg)
                - \beta D_{\text{KL}}(\pi_\theta || \pi_{\text{ref}})
            \Bigg\}
            \Bigg]
        \end{split}
        \end{equation}
        \item The importance ratio is formed \textbf{one token at a time}:
        \begin{equation}
            w_{i,t}(\theta) = \frac{\pi_\theta(o_{i,t} \mid q, o_{i,<t})}{\pi_{\theta_{\text{old}}}(o_{i,t} \mid q, o_{i,<t})}
        \end{equation}
    \framebreak
        \item KL divergence term:
        \begin{equation}
            D_{\text{KL}}(\pi_\theta || \pi_{\text{ref}}) =
            \frac{\pi_{\text{ref}}(o_{i,t} \mid q, o_{i,<t})}{\pi_\theta(o_{i,t} \mid q, o_{i,<t})} -
            \log \frac{\pi_{\text{ref}}(o_{i,t} \mid q, o_{i,<t})}{\pi_\theta(o_{i,t} \mid q, o_{i,<t})} - 1
        \end{equation}
        \item Under outcome supervision the checker fires once per response, so every token of $o_i$ inherits the same advantage, computed from the group rewards $\{r_1, r_2, \ldots, r_G\}$:
        \begin{equation}
            \hat{A}_{i,t} = \frac{r_i - \text{mean}(\{r_1, r_2, \ldots, r_G\})}{\text{std}(\{r_1, r_2, \ldots, r_G\})}
        \end{equation}
        \item $\epsilon$ and $\beta$ are hyper-parameters.
        \item The reward is granted once, for the whole response, while the ratio and the clip act on individual tokens. That mismatch is what GSPO attacks.
    \end{itemize}
\end{frame}

\begin{frame}{Why GRPO Is Popular for Open Reasoning Models}
  \begin{itemize}\itemsep0.3em
    \item[\textcolor{red}{$\blacktriangleright$}] \textbf{Fewer models in memory:} policy + reference (no critic, no RM when rules suffice).
    \item[\textcolor{red}{$\blacktriangleright$}] \textbf{Often faster to train} than PPO-style RLHF on comparable hardware.
    \item[\textcolor{red}{$\blacktriangleright$}] Powers DeepSeekMath, DeepSeek-R1-Zero, and many open reproductions.
    \item[\textcolor{red}{$\blacktriangleright$}] Cheap is relative: efficient kernels move an 8B long-context run from many GPUs to one node, while 4-bit QLoRA at roughly 1\,GB per billion parameters puts a 1.5B run in about 5\,GB.
  \end{itemize}
  \vspace{0.25em}
  \small
  \begin{tabular}{@{}lrr@{}}
    \toprule
    \textbf{Llama 3.1 8B, 20K context, 8 rollouts/prompt} & \textbf{Standard + FA2} & \textbf{Optimized} \\
    \midrule
    Training memory & 414\,GB & 42\,GB \\
    GRPO memory & 78.3\,GB & 9.8\,GB \\
    Inference + KV cache & 18.5\,GB & 2.5\,GB \\
    \midrule
    \textbf{Total VRAM} & \textbf{510.8\,GB} & \textbf{54.3\,GB} \\
    \bottomrule
  \end{tabular}
  \vspace{0.15em}

  {\scriptsize Source: Unsloth, \emph{Reinforcement Learning Guide}.}
\end{frame}

\begin{frame}{GRPO vs.\ PPO (Takeaways)}
  \small
  \setlength{\tabcolsep}{5pt}%
  \begin{tabular}{@{}p{0.22\textwidth}p{0.36\textwidth}p{0.36\textwidth}@{}}
    \toprule
    & \textbf{PPO (RLHF)} & \textbf{GRPO (RLVR)} \\
    \midrule
    Baseline & Critic / value net & Mean reward in a group \\
    Reward & Learned RM & Rules / verifiers \\
    Models in memory & Policy, ref., RM, critic & Policy, ref. (typical) \\
    Best for & Chat alignment & Math, code, logic \\
    \bottomrule
  \end{tabular}
  \vspace{0.5em}
  \begin{itemize}\itemsep0.3em
    \item[\textcolor{red}{$\blacktriangleright$}] GRPO does not replace DPO; DPO still wins when only preference pairs are available and there is no verifier.
  \end{itemize}
\end{frame}

\section{GSPO: Group Sequence Policy Optimization}

\begin{frame}{Motivation: Why GRPO Breaks at Scale}
    \begin{itemize}\itemsep0.3em
      \item[\textcolor{red}{$\blacktriangleright$}] Long RL runs on huge models need stable updates. Unstable RL ends in model collapse that reverting to an earlier checkpoint and retuning $\epsilon$ does not recover.
      \item[\textcolor{red}{$\blacktriangleright$}] The diagnosis (Zheng et al., 2025): importance sampling corrects a distribution mismatch only when you average over many samples from the behaviour policy. GRPO applies $w_{i,t}(\theta)$ at each token position using the \emph{single} token that was sampled there, so it corrects nothing and injects variance instead.
      \item[\textcolor{red}{$\blacktriangleright$}] That variance accumulates with response length and is amplified by per-token clipping, so the problem gets worse exactly as reasoning traces get longer.
      \item[\textcolor{red}{$\blacktriangleright$}] MoE models are worse again: on a 48-layer Qwen3-30B-A3B model, about 10\% of the experts activated for the same rollout differ after one gradient update, so token-level ratios swing hard. GRPO needs a workaround, Routing Replay, which caches the old policy's routing and replays it, to converge at all.
    \end{itemize}
\end{frame}

\begin{frame}{GSPO: Group Sequence Policy Optimization}
  \begin{itemize}\itemsep0.4em
    \item[\textcolor{red}{$\blacktriangleright$}] Rewards are usually sequence-level (correct answer, tests passed), so optimization should match that granularity.
    \item[\textcolor{red}{$\blacktriangleright$}] Same group-relative advantages as GRPO: compare $G$ completions on the same prompt.
    \item[\textcolor{red}{$\blacktriangleright$}] Design goals (\href{https://qwen.ai/blog?id=gspo}{Qwen blog}):
    \begin{itemize}\itemsep0.2em
      \item Stable, especially for MoE RL;
      \item Efficient, with better compute-to-gain than GRPO;
      \item Infrastructure-friendly, more robust to numeric precision issues in distributed RL.
    \end{itemize}
  \end{itemize}
\end{frame}

\begin{frame}[allowframebreaks]{GSPO: The Sequence-Level Objective}
  \begin{itemize}\itemsep0.3em
    \item[\textcolor{red}{$\blacktriangleright$}] One importance ratio per response, clipped once:
      \begin{equation}
        J_{\mathrm{GSPO}}(\theta) = \mathbb{E}_{x \sim \mathcal{D},\, \{y_i\}_{i=1}^G \sim \pi_{\theta_{\mathrm{old}}}(\cdot|x)}
        \left[ \frac{1}{G} \sum_{i=1}^{G} \min\big(s_i(\theta)\hat{A}_i,\; \mathrm{clip}(s_i(\theta), 1-\epsilon, 1+\epsilon)\hat{A}_i\big) \right]
      \end{equation}
    \item[\textcolor{red}{$\blacktriangleright$}] The ratio is the \textbf{length-normalized sequence likelihood ratio}:
      \begin{equation}
        s_i(\theta) = \left( \frac{\pi_\theta(y_i \mid x)}{\pi_{\theta_{\mathrm{old}}}(y_i \mid x)} \right)^{1/|y_i|}
        = \exp\!\left( \frac{1}{|y_i|} \sum_{t=1}^{|y_i|} \log \frac{\pi_\theta(y_{i,t} \mid x, y_{i,<t})}{\pi_{\theta_{\mathrm{old}}}(y_{i,t} \mid x, y_{i,<t})} \right)
      \end{equation}
    \framebreak
    \item[\textcolor{red}{$\blacktriangleright$}] $\hat{A}_i$ is the same group-relative advantage as GRPO. Every token of $y_i$ now carries the same gradient weight $s_i(\theta)$ rather than its own $w_{i,t}(\theta)$, which is where the variance goes.
    \item[\textcolor{red}{$\blacktriangleright$}] The $1/|y_i|$ exponent matters: without it the ratio swings with response length, and each length would need its own $\epsilon$.
    \item[\textcolor{red}{$\blacktriangleright$}] Because $s_i$ and $w_{i,t}$ live on different numerical scales, $\epsilon$ does too. On Qwen3-30B-A3B the paper clips GSPO at $3\times10^{-4}$ and $4\times10^{-4}$ (left, right) against $0.2$ and $0.27$ for GRPO.
  \end{itemize}
\end{frame}

\begin{frame}{GSPO: Algorithm Architecture (Training Loop)}
  \small
  \begin{enumerate}\itemsep0.35em
    \item \textbf{Sample} a batch of prompts $x \sim \mathcal{D}$.
    \item For each $x$, draw a \textbf{group} $\{y_i\}_{i=1}^G$ from the old policy $\pi_{\theta_{\mathrm{old}}}$.
    \item \textbf{Reward} each full response (verifier / rule checker) $\Rightarrow$ group-relative advantage $\hat{A}_i$ (same formula as GRPO).
    \item Build the sequence-level importance weight $s_i(\theta)$ from $\pi_\theta(y_i|x)$ vs.\ $\pi_{\theta_{\mathrm{old}}}(y_i|x)$.
    \item \textbf{Clip and optimize once per sequence} (not token-by-token); optional KL to $\pi_{\mathrm{ref}}$ as in PPO/GRPO stacks.
    \item \textbf{Variant:} GSPO-token reintroduces per-token advantages when they are needed (e.g.\ multi-turn RL), while keeping $s_i(\theta)$ as the clipping signal.
  \end{enumerate}
\end{frame}

\begin{frame}{GRPO vs.\ GSPO}
  \small
  \begin{tabular}{@{}p{0.26\textwidth}p{0.34\textwidth}p{0.34\textwidth}@{}}
    \toprule
    & \textbf{GRPO} & \textbf{GSPO} \\
    \midrule
    Importance ratio & Per-token $w_{i,t}(\theta)$ & Sequence-level $s_i(\theta)$ \\
    Clipping & Token-by-token & Once per sequence \\
    Advantage & Group-relative & Group-relative (same) \\
    Reward alignment & Mismatch on long CoT & Matches sequence rewards \\
    MoE RL & Needs Routing Replay & Stable in Qwen3-scale runs \\
    Typical use & DeepSeek-R1, dense RL & Qwen3, large MoE RL \\
    \bottomrule
  \end{tabular}
  \vspace{0.45em}
  \begin{itemize}\itemsep0.3em
    \item[\textcolor{red}{$\blacktriangleright$}] GSPO clips two orders of magnitude more tokens than GRPO and still trains more efficiently, which says the discarded token-level gradients were mostly noise.
  \end{itemize}
\end{frame}

\begin{frame}{After GSPO: Softer Gates, Other Aggregations}
  \begin{itemize}\itemsep0.4em
    \item[\textcolor{red}{$\blacktriangleright$}] GSPO's clip is still hard: a response contributes either its full gradient or none of it.
    \item[\textcolor{red}{$\blacktriangleright$}] \textbf{SAPO} (Soft Adaptive Policy Optimization, Gao et al., 2025) replaces the clip with a smooth, temperature-controlled gate, so an off-policy token is attenuated rather than the whole sequence discarded.
    \item[\textcolor{red}{$\blacktriangleright$}] \textbf{H\"olderPO} (Chen et al., 2026) reads GRPO and GSPO as two points in one family: aggregate the token ratios with the H\"older mean of order $p$, where $p=1$ (arithmetic) recovers GRPO and $p \to 0$ (geometric) recovers GSPO. Annealing $p$ downward over training trades early signal amplification for late variance control.
    \item[\textcolor{red}{$\blacktriangleright$}] The shared lesson: what settles these algorithms is where the reward is granted and how much variance the estimator can absorb, not the shape of the loss.
  \end{itemize}
\end{frame}
